\section{Discussion}\label{sec:discussion}

\subsection{The phrase is part of the label}

A category name is not a string, it is a distribution over the ways people refer to the
category, and an open-vocabulary detector is asked for the category by sampling that
distribution once. Our measurement puts a number on what that sampling costs: a third of
the score scale, a quarter of the binary decisions, and up to a third of the average
precision. This is not a defect of one checkpoint. It follows from the score being a
function of the text query, which is the defining property of the whole family.

The corollary for evaluation is worth stating separately, because it affects how results
in this literature should be read. An open-vocabulary detector is usually benchmarked with
one template, sometimes selected on the benchmark itself. If the choice of template moves
AP by \APGAP, then two systems compared under different templates are not comparable, and
a template tuned on the test split is a tuned hyper-parameter. Reporting the template, and
preferably the spread over several, costs one column and removes an ambiguity.

\subsection{What a calibrator can and cannot repair}

Proposition~\ref{prop:mono} draws a line that we found more useful than expected. Everything
the score \emph{says} --- the probability it assigns, the threshold at which it becomes a
decision, the expected number of mistakes in a review queue --- lives on the scale, and a
monotone map is free to move it. Everything the score \emph{orders} lives in the ranking,
and no monotone map can touch it. Two of our results are that line seen from either side:
PhC repairs the band by a factor of \FACTOR{} without moving AP at all, and no calibrator,
PhC included, can make a pooled threshold deliver its target on a wording whose ranking is
worse.

The two held-out experiments in Section~\ref{sec:exp-heldout} and
Section~\ref{sec:exp-transfer} put a second axis under the same distinction. A table indexed
by category generalises across wordings and not across vocabularies; a descriptor of the
phrase generalises across vocabularies and only partly across wordings. Neither relation
dominates the other, which is why we report both and claim neither as the general answer.

The practical lesson is that calibration and risk control are different jobs. Calibration
makes the number mean what it says; risk control makes the decision safe. A user who wants
the second cannot get it by improving the first, and the shortcut of tuning one threshold
on the mixed wordings and shipping it is exactly the case where the two diverge.

\subsection{Limitations}

\emph{Five wordings, and they are templates.} Our wordings are the ones that circulate as
prompt engineering, not free-form user language. Real queries are longer (``the dog on the
left'', ``a red car, side view'') and the descriptor we use would extrapolate to them
linearly, which we have not tested. The measurement protocol, however, applies unchanged:
whatever the phrasing, run it and compare.

\emph{One detector family.} Grounding DINO is a dual-encoder detector with a
text-conditioned decoder. A region--text model such as CLIP scoring a fixed proposal set
would respond differently in magnitude, and a model fine-tuned on captions differently
again. We expect the direction to survive, since the mechanism is shared, but we have not
measured it and do not claim it.

\emph{The descriptor needs the text embedding.} PhC reads the pooled last-layer text
embedding, which the detector computes anyway. A deployment behind an API that returns
boxes and scores only would need a proxy encoder for the phrase, and the calibration would
degrade by however much the proxy disagrees with the hidden encoder. This is the one place
where our method asks for something a black-box service might not expose.

\emph{The transfer run changes the vocabulary, not the domain.} LVIS is a different and much
larger label set over the same kind of photographs, and the 57 images of the split that carry
no annotation are the only ones missing from the run. It therefore tests the axis we care
about --- a category the calibrator has never seen --- and it does not test a change of image
domain; a detector deployed on microscopy or aerial imagery would face a shift that neither
LVIS nor COCO represents.

\emph{The guarantee is empirical across wordings.} Within a wording, conformalised
selection controls the false discovery rate under exchangeability of calibration and test
detections. Across wordings we do not have that assumption, and we do not claim a theorem:
Section~\ref{sec:exp-fdr} measures the gap. The worst-case threshold we recommend is valid
per wording, so it is covered by the existing guarantee; it is the \emph{pooled} threshold
that is not, and that is the one that fails.

\emph{The band is a choice.} We restrict attention to $p\in[0.3,0.9)$ because it is where
a reviewer spends their time, and we report global ECE alongside so the choice is visible.
The ranking of methods does not change between the two, but the magnitudes do.

\subsection{What we would recommend to someone shipping a detector}

Measure the spread first. Take the categories you care about, the wordings you expect, and
run the protocol of Section~\ref{sec:task} on a few hundred images: it costs one afternoon
and it tells you whether the problem is large in your domain, which it will be.

Then decide by what your deployment looks like. If the category set is small and fixed,
a per-category scale and bias is competitive with anything we tried, and cheap. If the
category set is open --- the reason to use an open-vocabulary detector at all --- the
per-category table has nothing to look up on the categories that matter, and a calibrator
that reads the phrase is the only one that transfers.

Finally, do not ship one threshold for all wordings. Select it per wording, and if you can
only ship one number, take the strictest of the per-wording thresholds rather than the
pooled one; the pooled threshold is optimistic by up to a factor of \FDROVERSHOOT{} on the
wording it was not tuned for, and the reviewer who types that wording will find out.

\section{Conclusion}

The score an open-vocabulary detector returns is a property of the phrase that requested
the category as much as of the region it describes. Changing only the phrasing moves the
score of the same object by \SPREADPTS{} on average, and the consequence is not cosmetic:
the threshold a user tuned no longer means what it meant, and a threshold selected on one
wording overshoots its false discovery target by up to \FDROVERSHOOT{} on another.

We proposed prompt-conditioned calibration, which predicts the two parameters of an
affine-in-logit map from a 33-dimensional descriptor of the phrase, built from the text
embedding the detector already computed. It has 68 parameters, needs no extra forward pass,
is strictly monotone within a run so that average precision is untouched by construction,
and reduces the error in the operating band by a factor of \FACTOR{} on COCO. Because its
parameters are indexed by a phrase rather than by a category, it transfers to a vocabulary
it has never seen, where a per-category table has nothing to look up.

We then separated calibration from risk control. A calibrated score makes the number mean
what it says; it does not make a shared threshold safe, because the false discovery rate at
a shared threshold is set by the order of the boxes and no monotone map can change that
order. A threshold chosen to satisfy the constraint inside every wording does hold, and
calibration pays for itself there, modestly: at the same target the calibrated score keeps
\WCGAIN{} more boxes than the uncalibrated one, at a smaller price in coverage. The gap between what calibration can do and what a
deployment needs is, we think, the more useful half of the result.
